\IfFileExists{stacks-project.cls}{%
\documentclass{stacks-project}
}{%
\documentclass{amsart}
}

% For dealing with references we use the comment environment
\usepackage{verbatim}
\newenvironment{reference}{\comment}{\endcomment}
%\newenvironment{reference}{}{}
\newenvironment{slogan}{\comment}{\endcomment}
\newenvironment{history}{\comment}{\endcomment}

% For commutative diagrams we use Xy-pic
\usepackage[all]{xy}

% We use 2cell for 2-commutative diagrams.
\xyoption{2cell}
\UseAllTwocells

% We use multicol for the list of chapters between chapters
\usepackage{multicol}

% This is generally recommended for better output
\usepackage{lmodern}
\usepackage[T1]{fontenc}

% For cross-file-references
\usepackage{xr-hyper}

% Package for hypertext links:
\usepackage{hyperref}

% For any local file, say "hello.tex" you want to link to please
% use \externaldocument[hello-]{hello}
\externaldocument[introduction-]{introduction}
\externaldocument[conventions-]{conventions}
\externaldocument[sets-]{sets}
\externaldocument[categories-]{categories}
\externaldocument[topology-]{topology}
\externaldocument[sheaves-]{sheaves}
\externaldocument[sites-]{sites}
\externaldocument[stacks-]{stacks}
\externaldocument[fields-]{fields}
\externaldocument[algebra-]{algebra}
\externaldocument[brauer-]{brauer}
\externaldocument[homology-]{homology}
\externaldocument[derived-]{derived}
\externaldocument[simplicial-]{simplicial}
\externaldocument[more-algebra-]{more-algebra}
\externaldocument[smoothing-]{smoothing}
\externaldocument[modules-]{modules}
\externaldocument[sites-modules-]{sites-modules}
\externaldocument[injectives-]{injectives}
\externaldocument[cohomology-]{cohomology}
\externaldocument[sites-cohomology-]{sites-cohomology}
\externaldocument[dga-]{dga}
\externaldocument[dpa-]{dpa}
\externaldocument[sdga-]{sdga}
\externaldocument[hypercovering-]{hypercovering}
\externaldocument[schemes-]{schemes}
\externaldocument[constructions-]{constructions}
\externaldocument[properties-]{properties}
\externaldocument[morphisms-]{morphisms}
\externaldocument[coherent-]{coherent}
\externaldocument[divisors-]{divisors}
\externaldocument[limits-]{limits}
\externaldocument[varieties-]{varieties}
\externaldocument[topologies-]{topologies}
\externaldocument[descent-]{descent}
\externaldocument[perfect-]{perfect}
\externaldocument[more-morphisms-]{more-morphisms}
\externaldocument[flat-]{flat}
\externaldocument[groupoids-]{groupoids}
\externaldocument[more-groupoids-]{more-groupoids}
\externaldocument[etale-]{etale}
\externaldocument[chow-]{chow}
\externaldocument[intersection-]{intersection}
\externaldocument[pic-]{pic}
\externaldocument[weil-]{weil}
\externaldocument[adequate-]{adequate}
\externaldocument[dualizing-]{dualizing}
\externaldocument[duality-]{duality}
\externaldocument[discriminant-]{discriminant}
\externaldocument[derham-]{derham}
\externaldocument[local-cohomology-]{local-cohomology}
\externaldocument[algebraization-]{algebraization}
\externaldocument[curves-]{curves}
\externaldocument[resolve-]{resolve}
\externaldocument[models-]{models}
\externaldocument[functors-]{functors}
\externaldocument[equiv-]{equiv}
\externaldocument[pione-]{pione}
\externaldocument[etale-cohomology-]{etale-cohomology}
\externaldocument[proetale-]{proetale}
\externaldocument[relative-cycles-]{relative-cycles}
\externaldocument[more-etale-]{more-etale}
\externaldocument[trace-]{trace}
\externaldocument[crystalline-]{crystalline}
\externaldocument[spaces-]{spaces}
\externaldocument[spaces-properties-]{spaces-properties}
\externaldocument[spaces-morphisms-]{spaces-morphisms}
\externaldocument[decent-spaces-]{decent-spaces}
\externaldocument[spaces-cohomology-]{spaces-cohomology}
\externaldocument[spaces-limits-]{spaces-limits}
\externaldocument[spaces-divisors-]{spaces-divisors}
\externaldocument[spaces-over-fields-]{spaces-over-fields}
\externaldocument[spaces-topologies-]{spaces-topologies}
\externaldocument[spaces-descent-]{spaces-descent}
\externaldocument[spaces-perfect-]{spaces-perfect}
\externaldocument[spaces-more-morphisms-]{spaces-more-morphisms}
\externaldocument[spaces-flat-]{spaces-flat}
\externaldocument[spaces-groupoids-]{spaces-groupoids}
\externaldocument[spaces-more-groupoids-]{spaces-more-groupoids}
\externaldocument[bootstrap-]{bootstrap}
\externaldocument[spaces-pushouts-]{spaces-pushouts}
\externaldocument[spaces-chow-]{spaces-chow}
\externaldocument[groupoids-quotients-]{groupoids-quotients}
\externaldocument[spaces-more-cohomology-]{spaces-more-cohomology}
\externaldocument[spaces-simplicial-]{spaces-simplicial}
\externaldocument[spaces-duality-]{spaces-duality}
\externaldocument[formal-spaces-]{formal-spaces}
\externaldocument[restricted-]{restricted}
\externaldocument[spaces-resolve-]{spaces-resolve}
\externaldocument[formal-defos-]{formal-defos}
\externaldocument[defos-]{defos}
\externaldocument[cotangent-]{cotangent}
\externaldocument[examples-defos-]{examples-defos}
\externaldocument[algebraic-]{algebraic}
\externaldocument[examples-stacks-]{examples-stacks}
\externaldocument[stacks-sheaves-]{stacks-sheaves}
\externaldocument[criteria-]{criteria}
\externaldocument[artin-]{artin}
\externaldocument[quot-]{quot}
\externaldocument[stacks-properties-]{stacks-properties}
\externaldocument[stacks-morphisms-]{stacks-morphisms}
\externaldocument[stacks-limits-]{stacks-limits}
\externaldocument[stacks-cohomology-]{stacks-cohomology}
\externaldocument[stacks-perfect-]{stacks-perfect}
\externaldocument[stacks-introduction-]{stacks-introduction}
\externaldocument[stacks-more-morphisms-]{stacks-more-morphisms}
\externaldocument[stacks-geometry-]{stacks-geometry}
\externaldocument[moduli-]{moduli}
\externaldocument[moduli-curves-]{moduli-curves}
\externaldocument[examples-]{examples}
\externaldocument[exercises-]{exercises}
\externaldocument[guide-]{guide}
\externaldocument[desirables-]{desirables}
\externaldocument[coding-]{coding}
\externaldocument[obsolete-]{obsolete}
\externaldocument[fdl-]{fdl}
\externaldocument[index-]{index}

% Theorem environments.
%
\theoremstyle{plain}
\newtheorem{theorem}[subsection]{Theorem}
\newtheorem{proposition}[subsection]{Proposition}
\newtheorem{lemma}[subsection]{Lemma}

\theoremstyle{definition}
\newtheorem{definition}[subsection]{Definition}
\newtheorem{example}[subsection]{Example}
\newtheorem{exercise}[subsection]{Exercise}
\newtheorem{situation}[subsection]{Situation}

\theoremstyle{remark}
\newtheorem{remark}[subsection]{Remark}
\newtheorem{remarks}[subsection]{Remarks}

\numberwithin{equation}{subsection}

% Macros
%
\def\lim{\mathop{\mathrm{lim}}\nolimits}
\def\colim{\mathop{\mathrm{colim}}\nolimits}
\def\Spec{\mathop{\mathrm{Spec}}}
\def\Hom{\mathop{\mathrm{Hom}}\nolimits}
\def\Ext{\mathop{\mathrm{Ext}}\nolimits}
\def\SheafHom{\mathop{\mathcal{H}\!\mathit{om}}\nolimits}
\def\SheafExt{\mathop{\mathcal{E}\!\mathit{xt}}\nolimits}
\def\Sch{\mathit{Sch}}
\def\Mor{\mathop{\mathrm{Mor}}\nolimits}
\def\Ob{\mathop{\mathrm{Ob}}\nolimits}
\def\Sh{\mathop{\mathit{Sh}}\nolimits}
\def\NL{\mathop{N\!L}\nolimits}
\def\CH{\mathop{\mathrm{CH}}\nolimits}
\def\proetale{{pro\text{-}\acute{e}tale}}
\def\etale{{\acute{e}tale}}
\def\QCoh{\mathit{QCoh}}
\def\Ker{\mathop{\mathrm{Ker}}}
\def\Im{\mathop{\mathrm{Im}}}
\def\Coker{\mathop{\mathrm{Coker}}}
\def\Coim{\mathop{\mathrm{Coim}}}

% Boxtimes
%
\DeclareMathSymbol{\boxtimes}{\mathbin}{AMSa}{"02}

%
% Macros for moduli stacks/spaces
%
\def\QCohstack{\mathcal{QC}\!\mathit{oh}}
\def\Cohstack{\mathcal{C}\!\mathit{oh}}
\def\Spacesstack{\mathcal{S}\!\mathit{paces}}
\def\Quotfunctor{\mathrm{Quot}}
\def\Hilbfunctor{\mathrm{Hilb}}
\def\Curvesstack{\mathcal{C}\!\mathit{urves}}
\def\Polarizedstack{\mathcal{P}\!\mathit{olarized}}
\def\Complexesstack{\mathcal{C}\!\mathit{omplexes}}
% \Pic is the operator that assigns to X its picard group, usage \Pic(X)
% \Picardstack_{X/B} denotes the Picard stack of X over B
% \Picardfunctor_{X/B} denotes the Picard functor of X over B
\def\Pic{\mathop{\mathrm{Pic}}\nolimits}
\def\Picardstack{\mathcal{P}\!\mathit{ic}}
\def\Picardfunctor{\mathrm{Pic}}
\def\Deformationcategory{\mathcal{D}\!\mathit{ef}}
\usepackage{cleveref}
\crefname{originallemma}{original lemma}{original lemmas}
\Crefname{originallemma}{Original lemma}{Original lemmas}
\makeatletter
\newlabel{homology-lemma-long-exact-sequence-chain}{{0111}{source}{}{sourcepointer-1}{}}
\newlabel{homology-lemma-long-exact-sequence-chain@cref}{{[originallemma][0][]0111}{[1][0][]source}{}{}{}}
\makeatother
\crefname{originallemma}{original lemma}{original lemmas}
\Crefname{originallemma}{Original lemma}{Original lemmas}
\makeatletter
\newlabel{homology-lemma-map-homology-homotopy}{{0110}{source}{}{sourcepointer-2}{}}
\newlabel{homology-lemma-map-homology-homotopy@cref}{{[originallemma][0][]0110}{[1][0][]source}{}{}{}}
\makeatother
\crefname{originallemma}{original lemma}{original lemmas}
\Crefname{originallemma}{Original lemma}{Original lemmas}
\makeatletter
\newlabel{categories-lemma-functorial-limit}{{002L}{source}{}{sourcepointer-3}{}}
\newlabel{categories-lemma-functorial-limit@cref}{{[originallemma][0][]002L}{[1][0][]source}{}{}{}}
\makeatother
\crefname{originallemma}{original lemma}{original lemmas}
\Crefname{originallemma}{Original lemma}{Original lemmas}
\makeatletter
\newlabel{categories-lemma-functorial-colimit}{{002K}{source}{}{sourcepointer-4}{}}
\newlabel{categories-lemma-functorial-colimit@cref}{{[originallemma][0][]002K}{[1][0][]source}{}{}{}}
\makeatother
\title{A termwise surjective quasi-isomorphism of simplicial abelian groups is a trivial Kan fibration}
\author{The Stacks Project Authors}
\date{stacks-project:simplicial.tex:local-snapshot-20261012; GNU FDL 1.2 or later}
\begin{document}
\maketitle
\noindent\textit{Editorial note.} The counted statement is the author's own lemma environment labelled lemma-qis-simplicial-abelian-groups (source lines 6053-6059); the proof environment begins on the second line after it and is its own proof, closing at source line 6099.\par\medskip
\noindent\textit{Editorial note.} The statement and the proof invoke two author results of the same chapter by label; both are retained as context spans and the generator expands the context transitively (budget 800 lines for this card), while reference targets in other chapters that carry canonical project tags are auto-declared as non-load-bearing source pointers.\par\medskip
\noindent\textit{Editorial note.} The retained preamble is the chapter's own minimal preamble (sources/stacks-simplicial/preamble.tex, taken with its real line count), a verbatim excerpt of the native Xy-pic setup, the native theorem declarations with the mathematical macros, and the project's own no-op reference environments; no macro was invented or redefined.\par\medskip
\noindent\textit{Editorial scope.} One original result of the Stacks Project's Simplicial Methods chapter, counted once: a termwise surjective homomorphism of simplicial abelian groups inducing a quasi-isomorphism on associated complexes is a trivial Kan fibration.\par\medskip
\noindent\textit{Allowed standard prerequisites (editorial).}\begin{itemize}
\item Simplicial abelian groups and associated chain complexes. The associated chain complex s(X) of a simplicial abelian group, its normalised complex N(X), and the definition of a quasi-isomorphism between them.. It is the setting and the acyclicity input (source lines 6055-6057, 6092-6095).
\item Trivial Kan fibrations. The lifting property defining a trivial Kan fibration, expressed as the datum of the horn maps a and b (the author's definition-trivial-kan, retained as context).. It is what the proof has to establish (source lines 6062-6073).
\item Simplicial groups are Kan complexes. Every simplicial group, in particular every simplicial abelian group, is a Kan complex (the author's lemma-simplicial-group-kan, retained as context).. It supplies the horn filler x in degree n (source lines 6087-6088).
\item The long exact sequence of homology. The long exact homology sequence of a short exact sequence of complexes (Homology, Lemma homology-lemma-long-exact-sequence-chain, declared as a tagged source pointer).. It shows that the kernel of f has vanishing homology, so that the reduction to Y = 0 stays a quasi-isomorphism (source lines 6080-6084).
\end{itemize}\medskip
\ifdefined\hypertarget\hypertarget{sourcepointer-1}{}\fi
\noindent\textit{Source reference.} Original lemma 0111: sources/stacks-chow/tags.txt, source lines 1347--1347. Cross-chapter pointer whose target carries the canonical tag 0111 in the project tag list; declared so the reference resolves inside the assembled package, while every step the counted argument uses remains inside the retained spans.\par\medskip
\ifdefined\hypertarget\hypertarget{sourcepointer-2}{}\fi
\noindent\textit{Source reference.} Original lemma 0110: sources/stacks-chow/tags.txt, source lines 1346--1346. Cross-chapter pointer whose target carries the canonical tag 0110 in the project tag list; declared so the reference resolves inside the assembled package, while every step the counted argument uses remains inside the retained spans.\par\medskip
\ifdefined\hypertarget\hypertarget{sourcepointer-3}{}\fi
\noindent\textit{Source reference.} Original lemma 002L: sources/stacks-chow/tags.txt, source lines 102--102. Cross-chapter pointer whose target carries the canonical tag 002L in the project tag list; declared so the reference resolves inside the assembled package, while every step the counted argument uses remains inside the retained spans.\par\medskip
\ifdefined\hypertarget\hypertarget{sourcepointer-4}{}\fi
\noindent\textit{Source reference.} Original lemma 002K: sources/stacks-chow/tags.txt, source lines 101--101. Cross-chapter pointer whose target carries the canonical tag 002K in the project tag list; declared so the reference resolves inside the assembled package, while every step the counted argument uses remains inside the retained spans.\par\medskip
\begin{definition}
\label{definition-trivial-kan}
A map $X \to Y$ of simplicial sets is called a {\it trivial Kan fibration}
if $X_0 \to Y_0$ is surjective and for all $n \geq 1$ and any commutative
solid diagram
$$
\xymatrix{
\partial \Delta[n] \ar[r] \ar[d] & X \ar[d] \\
\Delta[n] \ar[r] \ar@{-->}[ru] & Y
}
$$
a dotted arrow exists making the diagram commute.
\end{definition}

\begin{lemma}
\label{lemma-simplicial-group-kan}
Let $X$ be a simplicial group. Then $X$ is a Kan complex.
\end{lemma}

\begin{proof}
The following proof is basically just a translation into English
of the proof in the reference mentioned above.
Using the terminology as explained in the introduction to this section,
suppose $f : \Lambda_k[n] \to X$ is a morphism from a horn. Set
$x_i = f(\sigma_i) \in X_{n - 1}$ for $i = 0, \ldots, \hat k, \ldots, n$.
This means that for $i < j$ we have $d_i x_j = d_{j - 1} x_i$
whenever $i, j \not = k$.
We have to find an $x \in X_n$ such that $x_i = d_ix$ for
$i = 0, \ldots, \hat k, \ldots, n$.

\medskip\noindent
We first prove there exists a $u \in X_n$ such that $d_i u = x_i$
for $i < k$. This is trivial for $k = 0$. If $k > 0$, one defines by
induction an element $u^r \in X_n$ such that $d_i u^r = x_i$ for
$0 \leq i \leq r$. Start with $u^0 = s_0x_0$. If $r < k - 1$,
we set
$$
y^r = s_{r + 1}((d_{r + 1}u^r)^{-1}x_{r + 1}),\quad
u^{r + 1} = u^r y^r.
$$
An easy calculation shows that $d_iy^r = 1$ (unit element of the group
$X_{n - 1}$) for $i \leq r$ and
$d_{r + 1}y^r = (d_{r + 1}u^r)^{-1}x_{r + 1}$. It follows that
$d_iu^{r + 1} = x_i$ for $i \leq r + 1$. Finally, take $u = u^{k - 1}$
to get $u$ as promised.

\medskip\noindent
Next we prove, by induction on the integer $r$, $0 \leq r \leq n - k$,
there exists a $x^r \in X_n$ such that
$$
d_i x^r = x_i\quad\text{for }i < k\text{ and }i > n - r.
$$
Start with $x^0 = u$ for $r = 0$. Having defined $x^r$ for $r \leq n - k - 1$
we set
$$
z^r = s_{n - r - 1}((d_{n - r}x^r)^{-1}x_{n - r}),\quad x^{r + 1} = x^rz^r
$$
A simple calculation, using the given relations, shows that
$d_iz^r = 1$ for $i < k$ and $i > n - r$ and that
$d_{n - r}(z^r) = (d_{n - r}x^r)^{-1}x_{n - r}$. It follows that
$d_ix^{r + 1} = x_i$ for $i < k$ and $i > n - r - 1$. Finally, we
take $x = x^{n - k}$ which finishes the proof.
\end{proof}

\begin{lemma}
\label{lemma-quasi-isomorphism}
Let $\mathcal{A}$ be an abelian category.
Let $V$ be a simplicial object of $\mathcal{A}$.
The canonical morphism of chain complexes
$N(V) \to s(V)$ is a quasi-isomorphism.
In other words, the complex $D(V)$ of Lemma
\ref{lemma-decompose-associated-complexes} is acyclic.
\end{lemma}

\begin{proof}
Note that the result holds for $K(A, k)$ for
any object $A$ and any $k \geq 0$, by Lemmas
\ref{lemma-homology-eilenberg-maclane} and \ref{lemma-N-K}.
Consider the hypothesis $IH_{n, m}$:
for all $V$ such that $V_j = 0$ for
$j \leq m$ and all $i \leq n$ the map
$N(V) \to s(V)$ induces an isomorphism
$H_i(N(V)) \to H_i(s(V))$.

\medskip\noindent
To start of the induction, note that $IH_{n, n}$
is trivially true, because in that case $N(V)_n = 0$
and $s(V)_n = 0$.

\medskip\noindent
Assume $IH_{n, m}$, with $m \leq n$.
Pick a simplicial object $V$ such that
$V_j = 0$ for $j < m$. By Lemma \ref{lemma-eilenberg-maclane-object}
and Definition \ref{definition-eilenberg-maclane}
we have $K(V_m, m) = i_{m!} \text{sk}_mV$.
By Lemma \ref{lemma-n-skeleton-abelian} the natural morphism
$$
K(V_m, m) = i_{m!} \text{sk}_mV \to V
$$
is injective. Thus we get a short exact sequence
$$
0 \to K(V_m, m) \to V \to W \to 0
$$
for some $W$ with $W_i = 0$ for $i = 0, \ldots, m$. This short exact sequence
induces a morphism of short exact sequence of associated complexes
$$
\xymatrix{
0 \ar[r] &
N(K(V_m, m)) \ar[r] \ar[d] &
N(V) \ar[r] \ar[d] &
N(W) \ar[r] \ar[d] &
0 \\
0 \ar[r] &
s(K(V_m, m)) \ar[r] &
s(V) \ar[r] &
s(W) \ar[r] &
0
}
$$
see Lemmas \ref{lemma-s-exact} and \ref{lemma-N-exact}.
Hence we deduce the result for $V$ from the result
on the ends.
\end{proof}

\section{Simplicial objects and chain complexes}
\label{section-complexes}

\begin{definition}
\label{definition-eilenberg-maclane}
Let $\mathcal{A}$ be an abelian category.
Let $A$ be an object of $\mathcal{A}$ and
let $k$ be an integer $\geq 0$.
The {\it Eilenberg-Maclane object $K(A, k)$}
is given by the object $K(A, k) = i_{k!}U$
which is described in
Lemma \ref{lemma-eilenberg-maclane-object} above.
\end{definition}

\begin{lemma}
\label{lemma-N-K}
Let $\mathcal{A}$ be an abelian category.
Let $A$ be an object of $\mathcal{A}$ and
let $k$ be an integer. We have
$N(K(A, k))_i = A$ if $i = k$ and
$0$ else.
\end{lemma}

\begin{proof}
It is clear that $N(K(A, k))_i = 0$ when $i < k$
because $K(A, k)_i = 0$ in that case.
It is clear that $N(K(A, k))_k = A$ since
$K(A, k)_{k - 1} = 0$ and $K(A, k)_k = A$.
For $i > k$ we have $N(K(A, k))_i = 0$
by Lemma \ref{lemma-imshriek-abelian} and
the definition of $K(A, k)$, see Definition
\ref{definition-eilenberg-maclane}.
\end{proof}

\begin{lemma}
\label{lemma-N-exact}
The functor $N$ is exact.
\end{lemma}

\begin{proof}
By Lemma \ref{lemma-s-exact} and the functorial
decomposition of Lemma \ref{lemma-decompose-associated-complexes}.
\end{proof}

\begin{lemma}
\label{lemma-decompose-associated-complexes}
Let $\mathcal{A}$ be an abelian category.
Let $U$ be a simplicial object of $\mathcal{A}$.
The canonical morphism of chain complexes
$N(U) \to s(U)$ is split. In fact,
$$
s(U) = N(U) \oplus D(U)
$$
for some complex $D(U)$. The construction $U \mapsto D(U)$
is functorial.
\end{lemma}

\begin{proof}
Define $D(U)_n$ to be the image of
$$
\bigoplus\nolimits_{\varphi : [n] \to [m]\text{ surjective}, \ m < n} N(U_m)
\xrightarrow{\bigoplus U(\varphi)}
U_n
$$
which is a subobject of $U_n$ complementary to
$N(U_n)$ according to Lemma \ref{lemma-splitting-abelian-category} and
Definition \ref{definition-split}. We show that
$D(U)$ is a subcomplex. Pick a surjective
map $\varphi : [n] \to [m]$ with $m < n$ and consider
the composition
$$
N(U_m) \xrightarrow{U(\varphi)} U_n \xrightarrow{d_n} U_{n - 1}.
$$
This composition is the sum of the maps
$$
N(U_m) \xrightarrow{U(\varphi \circ \delta^n_i)} U_{n - 1}
$$
with sign $(-1)^i$, $i = 0, \ldots, n$.

\medskip\noindent
First we will prove by ascending induction on $m$,
$0 \leq m < n - 1$ that all the maps $U(\varphi \circ \delta^n_i)$
map $N(U_m)$ into $D(U)_{n - 1}$. (The case $m = n - 1$ is treated below.)
Whenever the map $\varphi \circ \delta^n_i : [n - 1] \to [m]$
is surjective then the image of $N(U_m)$ under $U(\varphi \circ \delta^n_i)$
is contained in $D(U)_{n - 1}$ by definition.
If $\varphi \circ \delta^n_i : [n - 1] \to [m]$ is not surjective,
set $j = \varphi(i)$ and observe that $i$ is the unique
index whose image under $\varphi$ is $j$. We may write
$\varphi \circ \delta^n_i = \delta^m_j \circ \psi \circ \delta^n_i$
for some $\psi : [n - 1] \to [m - 1]$. Hence
$U(\varphi \circ \delta^n_i) = U(\psi \circ \delta^n_i) \circ d^m_j$ which
is zero on $N(U_m)$ unless $j = m$. If $j = m$, then
$d^m_m(N(U_m)) \subset N(U_{m - 1})$ and hence
$U(\varphi \circ \delta^n_i)(N(U_m)) \subset
U(\psi \circ \delta^n_i)(N(U_{m - 1}))$ and we win
by induction hypothesis.

\medskip\noindent
To finish proving that $D(U)$ is a subcomplex
we still have to deal with the composition
$$
N(U_m) \xrightarrow{U(\varphi)} U_n \xrightarrow{d_n} U_{n - 1}.
$$
in case $m =  n - 1$. In this case $\varphi = \sigma^{n - 1}_j$
for some $0 \leq j \leq n - 1$ and $U(\varphi) = s^{n - 1}_j$.
Thus the composition is given by the sum
$$
\sum (-1)^i d^n_i \circ s^{n - 1}_j
$$
Recall from Remark \ref{remark-relations} that
$d^n_j \circ s^{n - 1}_j = d^n_{j + 1} \circ s^{n - 1}_j = \text{id}$
and these drop out because the corresponding terms have opposite signs.
The map $d^n_n \circ s^{n - 1}_j$, if $j < n - 1$, is equal to
$s^{n - 2}_j \circ d^{n - 1}_{n - 1}$. Since
$d^{n - 1}_{n - 1}$ maps $N(U_{n - 1})$ into $N(U_{n - 2})$,
we see that the image $d^n_n ( s^{n - 1}_j (N(U_{n - 1}))$
is contained in $s^{n - 2}_j(N(U_{n - 2}))$ which
is contained in $D(U_{n - 1})$ by definition. For all
other combinations of $(i, j)$ we have
either $d^n_i \circ s^{n - 1}_j = s^{n - 2}_{j - 1} \circ d^{n - 1}_i$
(if $i < j$), or
$d^n_i \circ s^{n - 1}_j = s^{n - 2}_j \circ d^{n - 1}_{i - 1}$
(if $n > i > j + 1$) and in these cases the map is zero because
of the definition of $N(U_{n - 1})$.
\end{proof}

\begin{lemma}
\label{lemma-eilenberg-maclane-object}
With $A$, $k$ and $U$ as above, so $U_i = 0$, $i < k$ and $U_k = A$.
\begin{enumerate}
\item Given a $k$-truncated simplicial object $V$
we have
$$
\Mor(U, V)
=
\{ f : A \to V_k \mid d^k_i \circ f = 0, \ i = 0, \ldots, k \}
$$
and
$$
\Mor(V, U)
=
\{ f : V_k \to A \mid f \circ s^{k - 1}_i = 0, \ i = 0, \ldots, k - 1 \}.
$$
\item The object $i_{k!} U$ has $n$th term equal to
$\bigoplus_\alpha A$ where $\alpha$ runs over all
surjective morphisms $\alpha : [n] \to [k]$.
\item For any $\varphi : [m] \to [n]$ the map
$i_{k!} U(\varphi)$ is described as the mapping
$\bigoplus_\alpha A \to \bigoplus_{\alpha'} A$
which maps to component corresponding to $\alpha : [n] \to [k]$
to zero if $\alpha \circ \varphi$ is not surjective and
by the identity to the component corresponding to
$\alpha \circ \varphi$ if it is surjective.
\item The object $\text{cosk}_k U$ has $n$th term equal to
$\bigoplus_\beta A$, where $\beta$ runs over all
injective morphisms $\beta : [k] \to [n]$.
\item For any $\varphi : [m] \to [n]$ the map
$\text{cosk}_k U(\varphi)$ is described as the mapping
$\bigoplus_\beta A \to \bigoplus_{\beta'} A$
which maps to component corresponding to $\beta : [k] \to [n]$
to zero if $\beta$ does not factor through $\varphi$ and
by the identity to each of the components corresponding to
$\beta'$ such that $\beta = \varphi \circ \beta'$
if it does.
\item The canonical map
$
c : i_{k !} U \to \text{cosk}_k U
$
in degree $n$ has $(\alpha, \beta)$ coefficient $A \to A$
equal to zero if $\alpha \circ \beta$ is not the identity
and equal to $\text{id}_A$ if it is.
\item The canonical map
$
c : i_{k !} U \to \text{cosk}_k U
$
is injective.
\end{enumerate}
\end{lemma}

\begin{proof}
The proof of (1) is left to the reader.

\medskip\noindent
Let us take the rules of (2) and (3)
as the definition of a simplicial object, call it $\tilde U$.
We will show that it is an incarnation of $i_{k!}U$.
This will prove (2), (3) at the same time. We have to show
that given a morphism $f : U \to \text{sk}_kV$
there exists a unique morphism $\tilde f : \tilde U \to V$
which recovers $f$ upon taking the $k$-skeleton.
From (1) we see that $f$ corresponds with a morphism
$f_k : A \to V_k$ which maps into the kernel of
$d^k_i$ for all $i$. For any surjective $\alpha : [n] \to [k]$
we set $\tilde f_\alpha : A \to V_n$ equal to the composition
$\tilde f_\alpha = V(\alpha) \circ f_k : A \to V_n$. We define
$\tilde f_n : \tilde U_n \to V_n$ as the sum of
the $\tilde f_\alpha$ over $\alpha : [n] \to [k]$ surjective.
Such a collection of $\tilde f_\alpha$ defines a morphism
of simplicial objects if and only if
for any $\varphi : [m] \to [n]$ the diagram
$$
\xymatrix{
\bigoplus_{\alpha : [n] \to [k]\text{ surjective}} A
\ar[r]_-{\tilde f_n}
\ar[d]_{(3)} &
V_n \ar[d]^{V(\varphi)} \\
\bigoplus_{\alpha' : [m] \to [k]\text{ surjective}} A
\ar[r]^-{\tilde f_m} &
V_m
}
$$
is commutative. Choosing $\varphi = \alpha$ shows our choice of
$\tilde f_\alpha$ is uniquely determined by $f_k$.
The commutativity in general may be checked for each summand
of the left upper corner separately. It is clear for the
summands corresponding to $\alpha$ where
$\alpha \circ \varphi$ is surjective, because those get
mapped by $\text{id}_A$ to the summand with
$\alpha' = \alpha \circ \varphi$, and we have
$\tilde f_{\alpha'} = V(\alpha') \circ f_k =
V(\alpha \circ \varphi) \circ f_k = V(\varphi) \circ \tilde f_\alpha$.
For those where $\alpha \circ \varphi$
is not surjective, we have to show that $V(\varphi) \circ \tilde f_\alpha = 0$.
By definition this is equal to
$V(\varphi) \circ V(\alpha) \circ f_k = V(\alpha \circ \varphi) \circ f_k$.
Since $\alpha \circ \varphi$ is not surjective we can write it
as $\delta^k_i \circ \psi$, and we deduce that
$V(\varphi) \circ V(\alpha) \circ f_k =
V(\psi) \circ d^k_i \circ f_k = 0$ see above.

\medskip\noindent
Let us take the rules of (4) and (5)
as the definition of a simplicial object, call it $\tilde U$.
We will show that it is an incarnation of $\text{cosk}_k U$.
This will prove (4), (5) at the same time. The argument is completely dual
to the proof of (2), (3) above, but we give it anyway.
We have to show
that given a morphism $f : \text{sk}_kV \to U$
there exists a unique morphism $\tilde f : V \to \tilde U$
which recovers $f$ upon taking the $k$-skeleton.
From (1) we see that $f$ corresponds with a morphism
$f_k : V_k \to A$ which is zero on the image of $s^{k - 1}_i$
for all $i$. For any injective $\beta : [k] \to [n]$
we set $\tilde f_\beta : V_n \to A$ equal to the composition
$\tilde f_\beta = f_k \circ V(\beta) : V_n \to A$. We define
$\tilde f_n : V_n \to \tilde U_n$ as the sum of
the $\tilde f_\beta$ over $\beta : [k] \to [n]$ injective.
Such a collection of $\tilde f_\beta$ defines a morphism
of simplicial objects if and only if
for any $\varphi : [m] \to [n]$ the diagram
$$
\xymatrix{
V_n
\ar[d]_{V(\varphi)}
\ar[r]_-{\tilde f_n}
&
\bigoplus_{\beta : [k] \to [n]\text{ injective}} A
\ar[d]^{(5)}
\\
V_m
\ar[r]^-{\tilde f_m}
&
\bigoplus_{\beta' : [k] \to [m]\text{ injective}} A
}
$$
is commutative. Choosing $\varphi = \beta$ shows our choice of
$\tilde f_\beta$ is uniquely determined by $f_k$.
The commutativity in general may be checked for each summand
of the right lower corner separately. It is clear for the
summands corresponding to $\beta'$ where
$\varphi \circ \beta'$ is injective, because these summands
get mapped into by exactly the summand with
$\beta = \varphi \circ \beta'$ and we have in that case
$\tilde f_{\beta'} \circ V(\varphi) =
f_k \circ V(\beta') \circ V(\varphi) =
f_k \circ V(\beta) = \tilde f_\beta$. For those where
$\varphi \circ \beta'$ is not injective,
we have to show that $\tilde f_{\beta'} \circ V(\varphi) = 0$.
By definition this is equal to
$f_k \circ V(\beta') \circ V(\varphi) =
f_k \circ V(\varphi \circ \beta')$.
Since $\varphi \circ \beta'$ is not injective we can write it
as $\psi \circ \sigma^{k - 1}_i$, and we deduce that
$f_k \circ V(\beta') \circ V(\varphi) =
f_k \circ s^{k - 1}_i \circ V(\psi) = 0$ see above.

\medskip\noindent
The composition $i_{k!}U \to \text{cosk}_kU$ is the
unique map of simplicial objects which is
the identity on $A = U_k = (i_{k!}U)_k = (\text{cosk}_kU)_k$.
Hence it suffices to check that the proposed rule defines
a morphism of simplicial objects.
To see this we have to show that
for any $\varphi : [m] \to [n]$ the diagram
$$
\xymatrix{
\bigoplus_{\alpha : [n] \to [k]\text{ surjective}} A
\ar[d]_{(3)}
\ar[r]_{(6)}
&
\bigoplus_{\beta : [k] \to [n]\text{ injective}} A
\ar[d]^{(5)}
\\
\bigoplus_{\alpha' : [m] \to [k]\text{ surjective}} A
\ar[r]^{(6)}
&
\bigoplus_{\beta' : [k] \to [m]\text{ injective}} A
}
$$
is commutative. Now we can think of this in terms of
matrices filled with only $0$'s and $1$'s as follows:
The matrix of (3) has a nonzero
$(\alpha', \alpha)$ entry if and only if
$\alpha' = \alpha \circ \varphi$. Likewise
the matrix of (5) has a nonzero
$(\beta', \beta)$ entry if and only if
$\beta = \varphi \circ \beta'$. The upper matrix
of (6) has a nonzero $(\alpha, \beta)$ entry if and only if
$\alpha \circ \beta = \text{id}_{[k]}$. Similarly for the
lower matrix of (6). The commutativity of the
diagram then comes down to computing the
$(\alpha, \beta')$ entry for both compositions
and seeing they are equal. This comes down to the
following equality
$$
\# \left\{
\beta
\mid
\beta = \varphi \circ \beta'
\text{ and }
\alpha \circ \beta = \text{id}_{[k]}
\right\}
=
\# \left\{
\alpha'
\mid
\alpha' = \alpha \circ \varphi
\text{ and }
\alpha' \circ \beta' = \text{id}_{[k]}
\right\}
$$
whose proof may safely be left to the reader.

\medskip\noindent
Finally, we prove (7). This follows directly from
Lemmas \ref{lemma-injective-map-simplicial-abelian},
\ref{lemma-recover-cosk}, \ref{lemma-recovering-U-for-real}
and \ref{lemma-imshriek-abelian}.
\end{proof}

\begin{lemma}
\label{lemma-homology-eilenberg-maclane}
Let $\mathcal{A}$ be an abelian category.
Let $A$ be an object of $\mathcal{A}$ and
let $k$ be an integer. We have
$H_i(s(K(A, k))) = A$ if $i = k$ and
$0$ else.
\end{lemma}

\begin{proof}
First, let us prove this if $k = 0$.
In this case we have $K(A, 0)_n = A$ for all $n$.
Furthermore, all the maps in this simplicial abelian
group are $\text{id}_A$, in other words $K(A, 0)$
is the constant simplicial object with value $A$.
The boundary maps $d_n = \sum_{i = 0}^n (-1)^i \text{id}_A
= 0$ if $n$ odd and $ = \text{id}_A$ if $n$ is even.
Thus $s(K(A, 0))$ looks like this
$$
\ldots \to A \xrightarrow{0} A \xrightarrow{1} A \xrightarrow{0} A \to 0
$$
and the result is clear.

\medskip\noindent
Next, we prove the result for all $k$ by induction.
Given the result for $k$ consider the short exact sequence
$$
0 \to K(A, k) \to E \to K(A, k + 1) \to 0
$$
from Lemma \ref{lemma-extension}.
By Lemma \ref{lemma-abelian} the associated sequence of
chain complexes is exact.
By Lemma \ref{lemma-homology-extension} we see that
$s(E)$ is acyclic. Hence the result for $k + 1$
follows from the long exact sequence of homology,
see Homology, Lemma \ref{homology-lemma-long-exact-sequence-chain}.
\end{proof}

\begin{lemma}
\label{lemma-n-skeleton-abelian}
Let $\mathcal{A}$ be an abelian category.
Let $U$ be a simplicial object of $\mathcal{A}$.
Let $n \geq 0$ be an integer.
The morphism $i_{n!} \text{sk}_n U \to U$ identifies
$i_{n!} \text{sk}_n U$ with the simplicial subobject
$U' \subset U$ defined in Lemma \ref{lemma-simplicial-abelian-n-skel-sub}.
\end{lemma}

\begin{proof}
By Lemma \ref{lemma-imshriek-abelian}
we have $N(i_{n!} \text{sk}_n U)_i = 0$
for $i > n$. The map
$i_{n!} \text{sk}_n U \to U$ is an isomorphism
in degrees $\leq n$, see Lemma \ref{lemma-recovering-U-for-real}.
Combined we conclude that the map $i_{n!} \text{sk}_n U \to U$
induces injective maps $N(i_{n!} \text{sk}_n U)_i \to N(U)_i$
for all $i$. Hence Lemma \ref{lemma-injective-map-simplicial-abelian}
applies. Thus $i_{n!} \text{sk}_n U \to U$
is injective. The result follows easily from this.
\end{proof}

\begin{lemma}
\label{lemma-s-exact}
The functor $s$ is exact.
\end{lemma}

\begin{proof}
Clear from Lemma \ref{lemma-abelian}.
\end{proof}

\begin{lemma}
\label{lemma-imshriek-abelian}
Let $\mathcal{A}$ be an abelian category
Let $U$ be an $m$-truncated simplicial object of
$\mathcal{A}$. For $n > m$ we have $N(i_{m!}U)_n = 0$.
\end{lemma}

\begin{proof}
Write $V = i_{m!}U$. Let $V' \subset V$ be the
simplicial subobject of $V$ with $V'_i = V_i$ for $i \leq m$
and $N(V'_i) = 0$ for $i > m$,
see Lemma \ref{lemma-simplicial-abelian-n-skel-sub}.
By the adjunction formula,
since $\text{sk}_m V' = U$, there is an inverse to the
injection $V' \to V$. Hence $V' = V$.
\end{proof}

\begin{definition}
\label{definition-split}
Let $\mathcal{C}$ be a category which admits finite nonempty coproducts.
We say a simplicial object $U$ of $\mathcal{C}$ is {\it split}
if there exist subobjects $N(U_m)$ of $U_m$, $m \geq 0$
with the property that
\begin{equation}
\label{equation-splitting}
\coprod\nolimits_{\varphi : [n] \to [m]\text{ surjective}}
N(U_m)
\longrightarrow
U_n
\end{equation}
is an isomorphism for all $n \geq 0$. If $U$ is an $r$-truncated
simplicial object of $\mathcal{C}$ then we say $U$ is {\it split}
if there exist subobjects $N(U_m)$ of $U_m$, $r \geq m \geq 0$
with the property that (\ref{equation-splitting})
is an isomorphism for $r \geq n \geq 0$.
\end{definition}

\begin{lemma}
\label{lemma-splitting-abelian-category}
Let $\mathcal{A}$ be an abelian category.
Let $U$ be a simplicial object in $\mathcal{A}$.
Then $U$ has a splitting obtained by taking $N(U_0) = U_0$ and
for $m \geq 1$ taking
$$
N(U_m) = \bigcap\nolimits_{i = 0}^{m - 1} \Ker(d^m_i).
$$
Moreover, this splitting is functorial on the category of
simplicial objects of $\mathcal{A}$.
\end{lemma}

\begin{proof}
For any object $A$ of $\mathcal{A}$ we obtain
a simplicial abelian group $\Mor_\mathcal{A}(A, U)$.
Each of these are canonically split by Lemma
\ref{lemma-splitting-simplicial-groups}. Moreover,
$$
N(\Mor_\mathcal{A}(A, U_m)) =
\bigcap\nolimits_{i = 0}^{m - 1} \Ker(d^m_i) =
\Mor_\mathcal{A}(A, N(U_m)).
$$
Hence we see that the morphism (\ref{equation-splitting})
becomes an isomorphism after applying the functor
$\Mor_\mathcal{A}(A, -)$ for any object of $\mathcal{A}$.
Hence it is an isomorphism by the Yoneda lemma.
\end{proof}

\begin{remark}
\label{remark-relations}
By abuse of notation we sometimes write $d_i : U_n \to U_{n - 1}$
instead of $d^n_i$, and similarly for $s_i : U_n \to U_{n + 1}$.
The relations among the morphisms $d^n_i$ and $s^n_i$
may be expressed as follows:
\begin{enumerate}
\item If $i < j$, then $d_i \circ d_j = d_{j - 1} \circ d_i$.
\item If $i < j$, then $d_i \circ s_j = s_{j - 1} \circ d_i$.
\item We have $\text{id} = d_j \circ s_j = d_{j + 1} \circ s_j$.
\item If $i > j + 1$, then $d_i \circ s_j = s_j \circ d_{i - 1}$.
\item If $i \leq j$, then $s_i \circ s_j = s_{j + 1} \circ s_i$.
\end{enumerate}
This means that whenever the compositions on both the left and the
right are defined then the corresponding equality should hold.
\end{remark}

\begin{lemma}
\label{lemma-injective-map-simplicial-abelian}
\begin{slogan}
The Dold-Kan normalization functor reflects
injectivity, surjectivity, and isomorphy.
\end{slogan}
Let $\mathcal{A}$ be an abelian category.
Let $f : U \to V$ be a morphism of
simplicial objects of $\mathcal{A}$.
If the induced morphisms $N(f)_i : N(U)_i \to N(V)_i$
are injective for all $i$, then $f_i$ is
injective for all $i$. Same holds with ``injective'' replaced
with ``surjective'', or ``isomorphism''.
\end{lemma}

\begin{proof}
This is clear from Lemma \ref{lemma-splitting-abelian-category}
and the definition of a splitting.
\end{proof}

\begin{lemma}
\label{lemma-recover-cosk}
Let $\mathcal{C}$ be a category with finite limits.
Let $U$ be an $n$-truncated simplicial object of $\mathcal{C}$.
The morphism $\text{sk}_n \text{cosk}_n U \to U$
is an isomorphism.
\end{lemma}

\begin{proof}
Combine Lemmas \ref{lemma-existence-cosk} and \ref{lemma-trivial-cosk}.
\end{proof}

\begin{lemma}
\label{lemma-recovering-U-for-real}
Let $\mathcal{C}$ be a category which has finite colimits.
Let $U$ be an $m$-truncated simplicial object of $\mathcal{C}$.
The map $U \to \text{sk}_m i_{m!}U$ is an isomorphism.
\end{lemma}

\begin{proof}
Combine Lemmas \ref{lemma-left-adjoint-exists} and \ref{lemma-recovering-U}.
\end{proof}

\begin{lemma}
\label{lemma-abelian}
Let $\mathcal{A}$ be an abelian category.
\begin{enumerate}
\item The categories $\text{Simp}(\mathcal{A})$ and
$\text{CoSimp}(\mathcal{A})$ are abelian.
\item A morphism of (co)simplicial objects
$f : A \to B$ is injective
if and only if each $f_n : A_n \to B_n$ is injective.
\item A morphism of (co)simplicial objects
$f : A \to B$ is surjective
if and only if each $f_n : A_n \to B_n$ is surjective.
\item A sequence of (co)simplicial objects
$$
A \xrightarrow{f} B \xrightarrow{g} C
$$
is exact at $B$ if and only if each sequence
$$
A_i \xrightarrow{f_i} B_i \xrightarrow{g_i} C_i
$$
is exact at $B_i$.
\end{enumerate}
\end{lemma}

\begin{proof}
Pre-additivity is easy. A final object is
given by $U_n = 0$ in all degrees.
Existence of direct products we saw in
Lemmas \ref{lemma-product} and
\ref{lemma-product-cosimplicial-objects}.
Kernels and cokernels are obtained by taking
termwise kernels and cokernels.
\end{proof}

\begin{lemma}
\label{lemma-extension}
Let $\mathcal{A}$ be an abelian category.
Let $A$ be an object of $\mathcal{A}$ and
let $k$ be an integer $\geq 0$. Consider the
simplicial object $E$ defined by the following rules
\begin{enumerate}
\item $E_n = \bigoplus_\alpha A$, where the
sum is over $\alpha : [n] \to [k + 1]$ whose
image is either $[k]$ or $[k + 1]$.
\item Given $\varphi : [m] \to [n]$ the map
$E_n \to E_m$ maps the summand corresponding
to $\alpha$ via $\text{id}_A$ to the summand
corresponding to $\alpha \circ \varphi$,
provided $\Im(\alpha \circ \varphi)$
is equal to $[k]$ or $[k + 1]$.
\end{enumerate}
Then there exists a short exact sequence
$$
0 \to K(A, k) \to E \to K(A, k + 1) \to 0
$$
which is term by term split exact.
\end{lemma}

\begin{proof}
The maps $K(A, k)_n \to E_n$ resp.\ $E_n \to K(A, k + 1)_n$
are given by the inclusion of direct sums, resp.\ projection
of direct sums which is obvious from the inclusions of index
sets. It is clear that these are maps of simplicial
objects.
\end{proof}

\begin{lemma}
\label{lemma-homology-extension}
Let $\mathcal{A}$ be an abelian category.
Let $A$ be an object of $\mathcal{A}$ and
let $k$ be an integer. Let $E$ be the object
described in Lemma \ref{lemma-extension}.
Then the complex $s(E)$ is acyclic.
\end{lemma}

\begin{proof}
For a morphism $\alpha : [n] \to [k + 1]$
we define $\alpha' : [n + 1] \to [k + 1]$ to be
the map such that $\alpha'|_{[n]} = \alpha$ and
$\alpha'(n + 1) = k + 1$. Note that if the
image of $\alpha$ is $[k]$ or $[k + 1]$, then
the image of $\alpha'$ is $[k + 1]$.
Consider the family of
maps $h_n : E_n \to E_{n + 1}$ which maps
the summand corresponding to $\alpha$ to
the summand corresponding to $\alpha'$ via
the identity on $A$.
Let us compute $d_{n + 1} \circ h_n - h_{n - 1} \circ d_n$.
We will first do this in case the category $\mathcal{A}$ is
the category of abelian groups.
Let us use the notation $x_\alpha$ to indicate
the element $x \in A$ in the summand of $E_n$ corresponding
to the map $\alpha$ occurring in the index set.
Let us also adopt the convention that
$x_\alpha$ designates the zero element of $E_n$
whenever $\Im(\alpha)$ is not $[k]$ or $[k + 1]$.
With these conventions we see that
$$
d_{n + 1}(h_n(x_\alpha)) =
\sum\nolimits_{i = 0}^{n + 1} (-1)^i x_{\alpha' \circ \delta^{n + 1}_i}
$$
and
$$
h_{n - 1}(d_n(x_\alpha)) =
\sum\nolimits_{i = 0}^n (-1)^i x_{(\alpha \circ \delta_i^n)'}
$$
It is easy to see that
$\alpha' \circ \delta^{n + 1}_i = (\alpha \circ \delta_i^n)'$
for $i = 0, \ldots, n$. It is also easy to see that
$\alpha' \circ \delta^{n + 1}_{n + 1} = \alpha$. Thus we
see that
$$
(d_{n + 1} \circ h_n - h_{n - 1} \circ d_n)(x_\alpha)
=
(-1)^{n + 1} x_\alpha
$$
These identities continue to hold if $\mathcal{A}$ is any abelian
category because they hold in the simplicial abelian group
$[n] \mapsto \Hom(A, E_n)$; details left to the reader.
We conclude that the identity map on $E$ is
homotopic to zero, with homotopy given by the
system of maps $h'_n = (-1)^{n + 1}h_n : E_n \to E_{n + 1}$.
Hence we see that $E$ is acyclic, for
example by Homology, Lemma \ref{homology-lemma-map-homology-homotopy}.
\end{proof}

\begin{lemma}
\label{lemma-simplicial-abelian-n-skel-sub}
Let $\mathcal{A}$ be an abelian category.
Let $U$ be a simplicial object of $\mathcal{A}$.
Let $n \geq 0$ be an integer.
The rule
$$
U'_m = \sum\nolimits_{\varphi : [m] \to [i], \ i\leq n} \Im(U(\varphi))
$$
defines a sub simplicial object $U' \subset U$ with $U'_i = U_i$
for $i \leq n$.
Moreover, $N(U'_m) = 0$ for all $m > n$.
\end{lemma}

\begin{proof}
Pick $m$, $i \leq n$ and some $\varphi : [m] \to [i]$.
The image under $U(\psi)$ of $\Im(U(\varphi))$
for any $\psi : [m'] \to [m]$ is
equal to the image of $U(\varphi \circ \psi)$ and
$\varphi \circ \psi : [m'] \to [i]$.
Hence $U'$ is a simplicial object.
Pick $m > n$. We have to show $N(U'_m) = 0$.
By definition of $N(U_m)$ and $N(U'_m)$ we have
$N(U'_m) = U'_m \cap N(U_m)$ (intersection of subobjects).
Since $U$ is split by Lemma \ref{lemma-splitting-abelian-category},
it suffices to show that $U'_m$ is contained in the sum
$$
\sum\nolimits_{\varphi : [m] \to [m']\text{ surjective}, \ m' < m}
\Im(U(\varphi)|_{N(U_{m'})}).
$$
By the splitting each $U_{m'}$ is the sum of images of
$N(U_{m''})$ via $U(\psi)$ for surjective maps
$\psi : [m'] \to [m'']$. Hence the displayed sum above
is the same as
$$
\sum\nolimits_{\varphi : [m] \to [m']\text{ surjective}, \ m' < m}
\Im(U(\varphi)).
$$
Clearly $U'_m$ is contained in this by the simple fact that
any $\varphi : [m] \to [i]$, $i \leq n$ occurring in the definition
of $U'_m$ may be factored as
$[m] \to [m'] \to [i]$ with $[m] \to [m']$ surjective
and $m' < m$ as in the last displayed sum above.
\end{proof}

\begin{lemma}
\label{lemma-splitting-simplicial-groups}
Let $U$ be a simplicial abelian group.
Then $U$ has a splitting obtained by taking $N(U_0) = U_0$ and
for $m \geq 1$ taking
$$
N(U_m) = \bigcap\nolimits_{i = 0}^{m - 1} \Ker(d^m_i).
$$
Moreover, this splitting is functorial on the category
of simplicial abelian groups.
\end{lemma}

\begin{proof}
By induction on $n$ we will show that the choice of $N(U_m)$
in the lemma guarantees that (\ref{equation-splitting}) is
an isomorphism for $m \leq n$. This is clear for $n = 0$.
In the rest of this proof we are going to
drop the superscripts from the maps $d_i$ and $s_i$ in order
to improve readability. We will also repeatedly use the relations
from Remark \ref{remark-relations}.

\medskip\noindent
First we make a general remark.
For $0 \leq i \leq m$ and $z \in U_m$ we have
$d_i(s_i(z)) = z$. Hence we can write
any $x \in U_{m + 1}$ uniquely as
$x = x' + x''$ with $d_i(x') = 0$
and $x'' \in \Im(s_i)$
by taking $x' = (x - s_i(d_i(x)))$ and
$x'' = s_i(d_i(x))$. Moreover, the element
$z \in U_m$ such that $x'' = s_i(z)$
is unique because $s_i$ is injective.

\medskip\noindent
Here is a procedure for decomposing any $x \in U_{n + 1}$.
First, write $x = x_0 + s_0(z_0)$ with $d_0(x_0) = 0$.
Next, write $x_0 = x_1 + s_1(z_1)$ with
$d_1(x_1) = 0$. Continue like this to get
\begin{eqnarray*}
x & = & x_0 + s_0(z_0), \\
x_0 & = & x_1 + s_1(z_1), \\
x_1 & = & x_2 + s_2(z_2), \\
\ldots & \ldots & \ldots \\
x_{n - 1} & = & x_n + s_n(z_n)
\end{eqnarray*}
where $d_i(x_i) = 0$ for all $i = n, \ldots, 0$.
By our general remark above all of the $x_i$
and $z_i$ are determined uniquely by $x$.
We claim that
$x_i \in
\Ker(d_0) \cap
\Ker(d_1) \cap
\ldots \cap
\Ker(d_i)$
and
$z_i \in
\Ker(d_0) \cap
\ldots \cap
\Ker(d_{i - 1})$
for $i = n, \ldots, 0$.
Here and in the following
an empty intersection of kernels indicates
the whole space; i.e.,
the notation
$z_0 \in \Ker(d_0) \cap
\ldots \cap
\Ker(d_{i - 1})$
when $i = 0$ means $z_0 \in U_n$ with no restriction.

\medskip\noindent
We prove this by ascending induction on $i$.
It is clear for $i = 0$ by construction of $x_0$ and $z_0$.
Let us prove it for $0 < i \leq n$ assuming the result for $i - 1$.
First of all we have $d_i(x_i) = 0$ by construction.
So pick a $j$ with $0 \leq j < i$. We have
$d_j(x_{i - 1}) = 0$ by induction. Hence
$$
0 = d_j(x_{i - 1})
= d_j(x_i) + d_j(s_i(z_i))
= d_j(x_i) + s_{i - 1}(d_j(z_i)).
$$
The last equality by the relations of Remark \ref{remark-relations}.
These relations also imply that
$d_{i - 1}(d_j(x_i)) = d_j(d_i(x_i)) = 0$
because $d_i(x_i)= 0$ by construction.
Then the uniqueness in the general remark above shows the equality
$0 = x' + x'' = d_j(x_i) + s_{i - 1}(d_j(z_i))$
can only hold if both terms are zero. We conclude that
$d_j(x_i) = 0$ and by injectivity of $s_{i - 1}$ we also
conclude that $d_j(z_i) = 0$. This proves the claim.

\medskip\noindent
The claim implies we can uniquely write
$$
x = s_0(z_0) + s_1(z_1) + \ldots + s_n(z_n) + x_n
$$
with $x_n \in N(U_{n + 1})$ and
$z_i \in \Ker(d_0) \cap \ldots \cap \Ker(d_{i - 1})$.
We can reformulate this as saying that we have found a direct
sum decomposition
$$
U_{n + 1}
=
N(U_{n + 1})
\oplus
\bigoplus\nolimits_{i = 0}^{i = n}
s_i\Big(\Ker(d_0) \cap \ldots \cap \Ker(d_{i - 1})\Big)
$$
with the property that
$$
\Ker(d_0) \cap \ldots \cap \Ker(d_j)
=
N(U_{n + 1}) \oplus
\bigoplus\nolimits_{i = j + 1}^{i = n}
s_i\Big(\Ker(d_0) \cap \ldots \cap \Ker(d_{i - 1})\Big)
$$
for $j = 0, \ldots, n$.
The result follows from this statement as follows.
Each of the $z_i$ in the expression for $x$
can be written uniquely as
$$
z_i = s_i(z'_{i, i}) + \ldots + s_{n - 1}(z'_{i, n - 1}) + z_{i, 0}
$$
with $z_{i, 0} \in N(U_n)$ and
$z'_{i, j} \in \Ker(d_0) \cap \ldots \cap \Ker(d_{j - 1})$.
The first few steps in the decomposition of $z_i$ are zero because
$z_i$ already is in the kernel of $d_0, \ldots, d_i$.
This in turn uniquely gives
$$
x = x_n + s_0(z_{0, 0}) + s_1(z_{1, 0}) + \ldots + s_n(z_{n, 0}) +
\sum\nolimits_{0 \leq i \leq j \leq n - 1} s_i(s_j(z'_{i, j})).
$$
Continuing in this fashion we see that we in the end obtain
a decomposition of $x$ as a sum of terms
of the form
$$
s_{i_1} s_{i_2} \ldots s_{i_k} (z)
$$
with $0 \leq i_1 \leq i_2 \leq \ldots \leq i_k \leq n - k + 1$ and
$z \in N(U_{n + 1 - k})$. This is exactly the required
decomposition, because any surjective map $[n + 1] \to [n + 1 - k]$
can be uniquely expressed in the form
$$
\sigma^{n - k}_{i_k} \ldots \sigma^{n - 1}_{i_2} \sigma^n_{i_1}
$$
with $0 \leq i_1 \leq i_2 \leq \ldots \leq i_k \leq n - k + 1$.
\end{proof}

\begin{lemma}
\label{lemma-existence-cosk}
If the category $\mathcal{C}$ has finite limits, then
$\text{cosk}_m$ functors exist for all $m$. Moreover,
for any $m$-truncated simplicial object $U$ the
simplicial object $\text{cosk}_mU$ is described
by the formula
$$
(\text{cosk}_mU)_n = \lim_{(\Delta/[n])_{\leq m}^{opp}} U(n)
$$
and for $\varphi : [n] \to [n']$ the map
$\text{cosk}_mU(\varphi)$ comes from the
identification $U(n') \circ \overline{\varphi} = U(n)$ above
via Categories, Lemma \ref{categories-lemma-functorial-limit}.
\end{lemma}

\begin{proof}
During the proof of this lemma we denote $\text{cosk}_mU$ the
simplicial object with $(\text{cosk}_mU)_n$ equal to
$\lim_{(\Delta/[n])_{\leq m}^{opp}} U(n)$.
We will conclude at the end of the proof that it does
satisfy the required mapping property.

\medskip\noindent
Suppose that $V$ is a simplicial object.
A morphism $\gamma : V \to \text{cosk}_mU$ is given by a sequence
of morphisms $\gamma_n : V_n \to (\text{cosk}_mU)_n$.
By definition of a limit, this is given by a
collection of morphisms $\gamma(\alpha) : V_n \to U_k$
where $\alpha$ ranges over all $\alpha : [k] \to [n]$
with $k \leq m$. These morphisms then also satisfy
the rules that
$$
\xymatrix{
V_n \ar[r]_{\gamma(\alpha)} &  U_k \\
V_{n'} \ar[r]^{\gamma(\alpha')} \ar[u]^{V(\varphi)} & U_{k'} \ar[u]_{U(\psi)}
}
$$
are commutative, given any $0 \leq k, k' \leq m$, $0 \leq n, n'$
and any $\psi : [k] \to [k']$, $\varphi : [n] \to [n']$,
$\alpha : [k] \to [n]$ and $\alpha' : [k'] \to [n']$ in $\Delta$
such that $\varphi \circ \alpha = \alpha' \circ \psi$.
Taking $n = k = k'$, $\varphi = \alpha'$, and $\alpha = \psi = \text{id}_{[k]}$
we deduce that $\gamma(\alpha') = \gamma(\text{id}_{[k]}) \circ V(\alpha')$.
In other words, the morphisms $\gamma(\text{id}_{[k]})$, $k \leq m$
determine the morphism $\gamma$. And it is easy to see that these
morphisms form a morphism $\text{sk}_m V \to U$.

\medskip\noindent
Conversely, given a morphism $\gamma : \text{sk}_m V \to U$,
we obtain a family of morphisms $\gamma(\alpha)$
where $\alpha$ ranges over all $\alpha : [k] \to [n]$
with $k \leq m$ by setting $\gamma(\alpha) =
\gamma(\text{id}_{[k]}) \circ V(\alpha)$. These morphisms
satisfy all the displayed commutativity restraints pictured
above, and hence give rise to a morphism $V \to \text{cosk}_m U$.
\end{proof}

\begin{lemma}
\label{lemma-trivial-cosk}
Let $\mathcal{C}$ be a category.
Let $U$ be an $m$-truncated simplicial object of $\mathcal{C}$.
For $n \leq m$ the limit $\lim_{(\Delta/[n])_{\leq m}^{opp}} U(n)$
exists and is canonically isomorphic to $U_n$.
\end{lemma}

\begin{proof}
This is true because the category $(\Delta/[n])_{\leq m}$
has an final object in this case, namely the identity
map $[n] \to [n]$.
\end{proof}

\begin{lemma}
\label{lemma-left-adjoint-exists}
Let $\mathcal{C}$ be a category which has finite colimits.
The functors $i_{m!}$ exist for all $m$.
Let $U$ be an $m$-truncated simplicial object of $\mathcal{C}$.
The simplicial object $i_{m!}U$
is described by the formula
$$
(i_{m!}U)_n = \colim_{([n]/\Delta)_{\leq m}^{opp}} U(n)
$$
and for $\varphi : [n] \to [n']$ the map
$i_{m!}U(\varphi)$ comes from the
identification $U(n) \circ \underline{\varphi} = U(n')$ above
via Categories, Lemma \ref{categories-lemma-functorial-colimit}.
\end{lemma}

\begin{proof}
In this proof we denote $i_{m!}U$ the simplicial object
whose $n$th term is given by the displayed formula of the
lemma. We will show it satisfies the adjointness property.

\medskip\noindent
Let $V$ be a simplicial object of $\mathcal{C}$.
Let $\gamma : U \to \text{sk}_mV$ be given.
A morphism
$$
\colim_{([n]/\Delta)_{\leq m}^{opp}} U(n) \to T
$$
is given by a compatible system of morphisms
$f_\alpha : U_k \to T$ where $\alpha : [n] \to [k]$
with $k \leq m$. Certainly, we have such a system of
morphisms by taking the compositions
$$
U_k \xrightarrow{\gamma_k} V_k \xrightarrow{V(\alpha)} V_n.
$$
Hence we get an induced morphism $(i_{m!}U)_n \to V_n$.
We leave it to the reader to see that these form a
morphism of simplicial objects $\gamma' : i_{m!}U \to V$.

\medskip\noindent
Conversely, given a morphism $\gamma' : i_{m!}U \to V$ we obtain
a morphism $\gamma : U \to \text{sk}_m V$ by setting
$\gamma_i : U_i \to V_i$ equal to the composition
$$
U_i
\xrightarrow{\text{id}_{[i]}}
\colim_{([i]/\Delta)_{\leq m}^{opp}} U(i)
\xrightarrow{\gamma'_i}
V_i
$$
for $0 \leq i \leq n$. We leave it to the reader to see that
this is the inverse of the construction above.
\end{proof}

\begin{lemma}
\label{lemma-recovering-U}
Let $\mathcal{C}$ be a category.
Let $U$ be an $m$-truncated simplicial object of $\mathcal{C}$.
For any $n \leq m$ the colimit
$$
\colim_{([n]/\Delta)_{\leq m}^{opp}} U(n)
$$
exists and is equal to $U_n$.
\end{lemma}

\begin{proof}
This is so because the category $([n]/\Delta)_{\leq m}$
has an initial object, namely $\text{id} : [n] \to [n]$.
\end{proof}

\begin{lemma}
\label{lemma-product}
If $U$ and $V$ are simplicial objects in the category $\mathcal{C}$,
and if $U \times V$ exists, then we have
$$
\Mor(W, U \times V) =
\Mor(W, U) \times
\Mor(W, V)
$$
for any third simplicial object $W$ of $\mathcal{C}$.
\end{lemma}

\begin{proof}
Omitted.
\end{proof}

\begin{lemma}
\label{lemma-product-cosimplicial-objects}
If $U$ and $V$ are cosimplicial objects in the category $\mathcal{C}$,
and if $U \times V$ exists, then we have
$$
\Mor(W, U \times V) =
\Mor(W, U) \times
\Mor(W, V)
$$
for any third cosimplicial object $W$ of $\mathcal{C}$.
\end{lemma}

\begin{proof}
Omitted.
\end{proof}

\begin{lemma}
\label{lemma-qis-simplicial-abelian-groups}
Let $f : X \to Y$ be a homomorphism of simplicial abelian groups
which is termwise surjective and induces a quasi-isomorphism on
associated chain complexes. Then $f$ is a trivial Kan fibration of
simplicial sets.
\end{lemma}

\begin{proof}
Consider a commutative solid diagram
$$
\xymatrix{
\partial \Delta[n] \ar[r]_a \ar[d] & X \ar[d] \\
\Delta[n] \ar[r]^b \ar@{-->}[ru] & Y
}
$$
as in Definition \ref{definition-trivial-kan}. The map $a$ corresponds
to $x_0, \ldots, x_n \in X_{n - 1}$ satisfying $d_i x_j = d_{j - 1} x_i$
for $i < j$. The map $b$ corresponds to an element $y \in Y_n$ such
that $d_iy = f(x_i)$. Our task is to produce an $x \in X_n$ such that
$d_ix = x_i$ and $f(x) = y$.

\medskip\noindent
Since $f$ is termwise surjective we can find $x \in X_n$ with $f(x) = y$.
Replace $y$ by $0 = y - f(x)$ and $x_i$ by $x_i - d_ix$. Then we
see that we may assume $y = 0$. In particular $f(x_i) = 0$.
In other words, we can replace $X$ by $\Ker(f) \subset X$
and $Y$ by $0$. This works, because by
Homology, Lemma \ref{homology-lemma-long-exact-sequence-chain}
the homology of the chain complex associated to $\Ker(f)$
is zero and hence $\Ker(f) \to 0$ induces a quasi-isomorphism
on associated chain complexes.

\medskip\noindent
Since $X$ is a Kan complex (Lemma \ref{lemma-simplicial-group-kan})
we can find $x \in X_n$ with $d_i x = x_i$ for $i = 0, \ldots, n - 1$.
After replacing $x_i$ by $x_i - d_ix$ for $i = 0, \ldots, n$ we
may assume that $x_0 = x_1 = \ldots = x_{n - 1} = 0$.
In this case we see that $d_i x_n = 0$ for $i = 0, \ldots, n - 1$.
Thus $x_n \in N(X)_{n - 1}$ and lies in the kernel of the differential
$N(X)_{n - 1} \to N(X)_{n - 2}$. Here $N(X)$ is the normalized chain
complex associated to $X$, see Section \ref{section-complexes}.
Since $N(X)$ is quasi-isomorphic to $s(X)$
(Lemma \ref{lemma-quasi-isomorphism})
and thus acyclic we find $x \in N(X_n)$ whose differential is $x_n$.
This $x$ answers the question posed by the lemma and we are done.
\end{proof}

\end{document}
